\chapter{Formulación de Elementos Finitos de Barra}
\label{chap:cap06-formulacion-fem-barras}

\section{De la forma débil a la aproximación discreta}

En el Capítulo 5 se estableció que el Principio de los Trabajos Virtuales (PTV) es equivalente, en estática, a las ecuaciones diferenciales de equilibrio junto con las condiciones de borde naturales: si un campo de tensiones $\bm\sigma$ satisface
\[
\delta W_{\text{ext}} = \int_{\Gamma_\sigma} \bar t_i\,\delta u_i \, d\Gamma + \int_\Omega \rho b_i \,\delta u_i \, d\Omega + P_i\,\delta u_i
= \int_\Omega \sigma_{ij}\,\delta\varepsilon_{ij}\, d\Omega = \delta W_{\text{int}}
\]
para todo campo de desplazamientos virtuales $\delta\bm u$ cinemáticamente admisible, entonces $\bm\sigma$ está en equilibrio. Esta forma integral tiene una ventaja decisiva frente a la forma diferencial: solo exige que $\bm u$ sea derivable una vez (mientras que la ecuación de equilibrio puntual exige dos derivadas de $\bm u$), y por lo tanto admite soluciones aproximadas construidas con funciones más simples que las que requeriría resolver directamente la ecuación diferencial.

Sin embargo, el PTV tal como fue enunciado es una igualdad entre integrales sobre un dominio continuo $\Omega$, válida para un espacio de funciones de dimensión infinita ($\delta\bm u$ arbitrario). Para convertirlo en un problema que un computador pueda resolver, es necesario reducir ese espacio de funciones a uno de dimensión \emph{finita}. Esta reducción es el segundo ingrediente fundamental del MEF, ya anunciado al cierre del Capítulo 5: la \textbf{interpolación del campo de desplazamientos mediante funciones de forma}, en términos de un número finito de parámetros (los desplazamientos nodales).

Este capítulo desarrolla ese ingrediente en detalle para el caso particular del elemento finito de barra, deduciendo de manera sistemática la matriz de rigidez $\bm K^{(e)}$ y el vector de fuerzas nodales equivalentes $\bm F^{(e)}$ para los elementos de dos y de tres nodos, y generalizando luego a elementos de un número arbitrario de nodos mediante polinomios de Lagrange. Un resultado central que se verificará explícitamente es que, para el elemento lineal de dos nodos, la matriz de rigidez obtenida por esta vía \emph{coincide exactamente} con la deducida en el Capítulo 4 a partir de argumentos de equilibrio y compatibilidad directos sobre la barra --- una validación cruzada que confirma la consistencia entre ambos caminos (análisis matricial clásico y formulación variacional vía PTV).

\section{Interpolación del campo de desplazamientos}

Sea un elemento finito con $n$ nodos. La idea central de la interpolación consiste en aproximar el campo de desplazamientos $\bm u$ en cualquier punto interior del elemento como una combinación lineal de los desplazamientos nodales $\bm U_i$, ponderados por funciones de interpolación (o \emph{funciones de forma}) $N_i$ asociadas a cada nodo:
\[
u = N_1 U_1 + N_2 U_2 + \dots + N_n U_n = N_i U_i ,
\]
donde se ha escrito la expresión para una componente escalar de $\bm u$ (en el caso de la barra, la única componente no nula es el desplazamiento axial $w$). Las funciones $N_i$ son, en general, polinomios que dependen únicamente de la geometría del elemento y de la posición del nodo $i$ dentro de él; no dependen de la solución del problema. Esta es la generalización natural de la interpolación lineal ya usada implícitamente en el Capítulo 4 al describir el campo $w$ dentro de una barra en función de sus valores extremos $W_1^{(e)}$ y $W_2^{(e)}$.

\begin{ecnimportante}{Interpolación del campo de desplazamientos vía funciones de forma}{cap06-interpolacion-desplazamientos}
Para un elemento finito de $n$ nodos, el campo de desplazamientos en su interior se aproxima como
\[
\bm u = \bm N \, \bm U^{(e)}, \qquad
\bm N = \begin{bmatrix} N_1 & N_2 & \cdots & N_n \end{bmatrix},\qquad
\bm U^{(e)} = \begin{bmatrix} U_1 \\ U_2 \\ \vdots \\ U_n \end{bmatrix},
\]
donde $\bm N$ es la \emph{matriz de funciones de interpolación} (elemental) y $\bm U^{(e)}$ es el vector de desplazamientos nodales del elemento. Para el caso particular de la barra, con desplazamiento axial $w$:
\[
w = N_1^{(e)} W_1^{(e)} + N_2^{(e)} W_2^{(e)} + \dots + N_n^{(e)} W_n^{(e)} = \bm N^{(e)} \bm U^{(e)} .
\]
Las funciones de forma deben satisfacer la propiedad de \emph{interpolación nodal}
\[
N_i^{(e)}(z_j) = \delta_{ij},
\]
es decir, cada $N_i$ vale $1$ en su propio nodo y $0$ en todos los demás nodos del elemento --- condición que garantiza que $\bm U^{(e)}$ contiene, en efecto, los valores físicos del campo $w$ en cada nodo.
\end{ecnimportante}

A los elementos que solo garantizan continuidad del propio campo de desplazamientos entre elementos adyacentes (sin exigir continuidad de sus derivadas) se los denomina \emph{elementos de clase $C^0$}. Todos los elementos de barra desarrollados en este capítulo son de esta clase, lo cual será relevante en el Capítulo 7 al discutir los requisitos de convergencia.

\section{Matriz deformación-desplazamiento}

Una vez interpolado el campo de desplazamientos, es necesario obtener el campo de deformaciones asociado, ya que este es el que aparece en el trabajo virtual interno $\delta W_{\text{int}} = \int_\Omega \sigma_{ij}\,\delta \varepsilon_{ij}\, d\Omega$. En una barra sometida únicamente a esfuerzo axial, todas las componentes del tensor de deformación son nulas salvo $\varepsilon_{zz}$, que se obtiene derivando el campo interpolado:
\[
\varepsilon_{zz} = \pdv{w}{z} = \pdv{N_1^{(e)}}{z} W_1^{(e)} + \pdv{N_2^{(e)}}{z} W_2^{(e)} + \dots + \pdv{N_n^{(e)}}{z} W_n^{(e)} .
\]
Como cada $N_i^{(e)}$ es función únicamente de $z$, la derivada parcial se reduce en este caso a la derivada total; se conserva la notación $\partial/\partial z$ por coherencia con la formulación general del PTV en $\mathbb R^3$ desarrollada en el Capítulo 5, de la cual la barra es un caso particular unidimensional.

\begin{ecnimportante}{Matriz deformación-desplazamiento}{cap06-matriz-B}
Escribiendo la relación anterior en forma matricial,
\[
\varepsilon_{zz} = \begin{bmatrix} \dfrac{\partial N_1^{(e)}}{\partial z} & \dfrac{\partial N_2^{(e)}}{\partial z} & \cdots & \dfrac{\partial N_n^{(e)}}{\partial z} \end{bmatrix} \bm U^{(e)} = \bm B^{(e)} \bm U^{(e)},
\]
se define la \emph{matriz de deformación} (o matriz deformación-desplazamiento) del elemento como
\[
\bm B^{(e)} = \pdv{\bm N^{(e)}}{z} ,
\]
es decir, la matriz que se obtiene derivando término a término la matriz de funciones de forma respecto de la coordenada axial. El mismo operador aplicado al campo de desplazamientos virtuales $\delta w = \bm N^{(e)}\,\delta \bm U^{(e)}$ entrega la deformación virtual $\delta\varepsilon_{zz} = \bm B^{(e)}\, \delta \bm U^{(e)}$, consistente con la manera en que $\delta$ opera de forma análoga a un diferencial (Capítulo 5).
\end{ecnimportante}

\section{Matriz constitutiva reducida al caso de la barra}

El tercer ingrediente necesario para cerrar el problema es la relación constitutiva entre tensión y deformación. En la teoría general desarrollada en el Capítulo 3 para un medio continuo tridimensional, la ley de Hooke generalizada se escribe en forma matricial como $\bm\sigma = \bm C\,\bm\varepsilon$, con $\bm C$ la matriz constitutiva elástica (de $6\times 6$ en notación de Voigt para el caso general, o reducida según el estado de tensión o deformación que corresponda).

En una barra prismática sometida únicamente a carga axial se admite la hipótesis de estado uniaxial de tensión: las caras laterales están libres de tracción y, por lo tanto, todas las componentes de $\bm\sigma$ son nulas salvo $\sigma_{zz}$. La ley de Hooke se reduce entonces a un escalar,
\[
\sigma_{zz} = E\,\varepsilon_{zz} = \underbrace{E^{(e)}}_{\bm C^{(e)}} \, \bm B^{(e)}\, \bm U^{(e)} ,
\]
de modo que la \emph{matriz constitutiva elemental} de la barra es simplemente $\bm C^{(e)} = E^{(e)}$, el módulo de Young del material del elemento. Esta reducción --- de una matriz $6\times 6$ en el caso general a un escalar en el caso de la barra --- es la misma simplificación que ya se había usado implícitamente en el Capítulo 4 al escribir $\sigma_{zz} = E\varepsilon_{zz}$ para obtener la rigidez axial $K^{(e)} = E^{(e)}A^{(e)}/L^{(e)}$.

\section{De la forma débil a las matrices elementales: la fórmula maestra}
\label{sec:cap06-formula-maestra}

Con los tres ingredientes anteriores --- interpolación $\bm u = \bm N \bm U$, deformación $\bm\varepsilon = \bm B \bm U$ y constitutiva $\bm\sigma = \bm C \bm\varepsilon$ --- se puede ahora sustituir directamente en el PTV aplicado a un elemento $(e)$, suprimiendo por simplicidad el superíndice $(e)$ como se hizo en el Capítulo 5:
\[
\delta W_{\text{int}} = \delta W_{\text{ext}} \qquad \Longrightarrow \qquad
\int_\Omega \delta\bm\varepsilon^{\mathsf T} \bm\sigma \, d\Omega = \int_{\Gamma_\sigma} \delta\bm u^{\mathsf T} \bar{\bm t} \, d\Gamma + \int_\Omega \delta\bm u^{\mathsf T} \rho \bm b \, d\Omega + \delta \bm U^{\mathsf T} \bm P .
\]

\paragraph{Trabajo virtual interno.} Sustituyendo $\delta\bm\varepsilon = \bm B\,\delta\bm U$ y $\bm\sigma = \bm C\bm\varepsilon = \bm C\bm B\bm U$:
\[
\int_\Omega \delta\bm\varepsilon^{\mathsf T}\bm\sigma\, d\Omega
= \int_\Omega (\bm B\,\delta\bm U)^{\mathsf T} \bm C \bm B \bm U \, d\Omega
= \delta\bm U^{\mathsf T} \underbrace{\left(\int_\Omega \bm B^{\mathsf T}\bm C\bm B\, d\Omega\right)}_{\bm K} \bm U ,
\]
donde se usó que $\delta\bm U$ (vector de desplazamientos nodales virtuales) es constante dentro de la integral sobre $\Omega$, ya que no depende de la posición $z$, y se pudo sacar fuera de la integral junto con $\bm U$.

\paragraph{Trabajo virtual externo.} Análogamente, sustituyendo $\delta\bm u = \bm N\,\delta\bm U$:
\[
\int_{\Gamma_\sigma} \delta\bm u^{\mathsf T}\bar{\bm t}\, d\Gamma + \int_\Omega \delta\bm u^{\mathsf T} \rho\bm b\, d\Omega + \delta\bm U^{\mathsf T}\bm P
= \delta\bm U^{\mathsf T} \underbrace{\left( \int_{\Gamma_\sigma} \bm N^{\mathsf T} \bar{\bm t}\, d\Gamma + \int_\Omega \bm N^{\mathsf T} \rho \bm b\, d\Omega + \bm P \right)}_{\bm F} .
\]

\paragraph{Ecuación resultante.} Igualando ambos trabajos virtuales,
\[
\delta \bm U^{\mathsf T} \bm K \bm U = \delta\bm U^{\mathsf T}\bm F \qquad\Longrightarrow\qquad \delta\bm U^{\mathsf T}\left(\bm K\bm U - \bm F\right) = \bm 0 ,
\]
y como $\delta\bm U^{\mathsf T}$ es arbitrario (cualquier desplazamiento virtual cinemáticamente admisible), el término entre paréntesis debe anularse idénticamente, lo que entrega la ecuación de equilibrio elemental $\bm K\bm U = \bm F$ --- exactamente la misma estructura obtenida en el Capítulo 4 por argumentos de equilibrio directo, pero ahora derivada de manera sistemática y generalizable a cualquier tipo de elemento.

\begin{ecnimportante}{Matrices elementales de rigidez y de fuerzas nodales equivalentes}{cap06-formula-maestra}
Para cualquier elemento finito, las matrices que gobiernan su comportamiento discreto se obtienen de la sustitución de $\bm u = \bm N\bm U$, $\bm\varepsilon = \bm B\bm U$ en el PTV:
\[
\bm K = \int_\Omega \bm B^{\mathsf T}\bm C\bm B \, d\Omega, \qquad
\bm F = \int_{\Gamma_\sigma} \bm N^{\mathsf T} \bar{\bm t}\, d\Gamma + \int_\Omega \bm N^{\mathsf T} \rho \bm b\, d\Omega + \bm P .
\]
Particularizado al elemento de barra, con $\bar t_z$ la carga axial distribuida por unidad de longitud aplicada sobre el elemento, $b_z$ la fuerza de cuerpo por unidad de masa y $\bm P$ las cargas nodales concentradas:
\[
\bm K = \int_\Omega \bm B^{\mathsf T} E \bm B\, d\Omega, \qquad
\bm F = \int_{L} \bm N^{\mathsf T} \bar t_z\, dz + \int_\Omega \bm N^{\mathsf T} \rho\, b_z \, d\Omega + \bm P .
\]
Este par de fórmulas --- válido para \emph{cualquier} número de nodos y cualquier orden de interpolación --- es el resultado central de este capítulo: reduce el cálculo de las matrices elementales a la elección de $\bm N$ (y, por lo tanto, de $\bm B$) y a la evaluación de dos integrales.
\end{ecnimportante}

Es interesante notar la simetría de $\bm K$: dado que $\bm K = \int \bm B^{\mathsf T} C \bm B \,d\Omega$ y $C$ es un escalar (o, en el caso general, una matriz simétrica), $\bm K^{\mathsf T} = \int \bm B^{\mathsf T} C^{\mathsf T} \bm B\, d\Omega = \bm K$, propiedad ya observada en el análisis matricial del Capítulo 4 y que se preserva automáticamente en la formulación por PTV.

\section{Elemento lineal de dos nodos}
\label{sec:cap06-elemento-2nodos}

Considérese un elemento de barra con dos nodos, situados en $z=0$ (nodo 1) y $z=L^{(e)}$ (nodo 2), de módulo de Young $E^{(e)}$ y área de sección $A^{(e)}$ (Figura~\ref{fig:cap06-gdl-elementos}a). Con solo dos nodos, el campo de desplazamientos axiales admite, como máximo, una interpolación \emph{lineal}:
\[
w = a_0 + a_1 z .
\]
Los coeficientes $a_0, a_1$ se determinan imponiendo que $w$ tome los valores nodales prescritos:
\[
w\big|_{z=0} = W_1^{(e)} = a_0, \qquad w\big|_{z=L^{(e)}} = W_2^{(e)} = a_0 + a_1 L^{(e)} \quad\Longrightarrow\quad a_0 = W_1^{(e)}, \quad a_1 = \frac{W_2^{(e)} - W_1^{(e)}}{L^{(e)}} .
\]
Sustituyendo de vuelta y reagrupando en términos de $W_1^{(e)}$ y $W_2^{(e)}$:
\[
w = \left(1 - \frac{z}{L^{(e)}}\right) W_1^{(e)} + \frac{z}{L^{(e)}} W_2^{(e)} = \underbrace{\left(1-\frac{z}{L^{(e)}}\right)}_{N_1^{(e)}} W_1^{(e)} + \underbrace{\frac{z}{L^{(e)}}}_{N_2^{(e)}} W_2^{(e)} .
\]
Se verifica de inmediato que $N_1^{(e)}(0)=1$, $N_1^{(e)}(L^{(e)})=0$, $N_2^{(e)}(0)=0$, $N_2^{(e)}(L^{(e)})=1$, es decir, se cumple la propiedad de interpolación nodal $N_i^{(e)}(z_j)=\delta_{ij}$ exigida en la \cref{eq:cap06-interpolacion-desplazamientos}. Además, $N_1^{(e)}+N_2^{(e)}=1$ para todo $z$, propiedad que en el Capítulo 7 se mostrará que es equivalente a la capacidad del elemento de representar exactamente un movimiento de sólido rígido.

La matriz de deformación se obtiene derivando:
\[
\bm B^{(e)} = \begin{bmatrix} \dfrac{\partial N_1^{(e)}}{\partial z} & \dfrac{\partial N_2^{(e)}}{\partial z} \end{bmatrix} = \begin{bmatrix} -\dfrac{1}{L^{(e)}} & \dfrac{1}{L^{(e)}} \end{bmatrix} ,
\]
que es, notablemente, \emph{constante} dentro del elemento --- consecuencia directa de que $\bm N$ sea lineal. Esto implica que el elemento lineal de dos nodos representa un estado de deformación (y por lo tanto de tensión) \emph{uniforme} en su interior, aproximación que solo será exacta si el campo de deformaciones real es efectivamente constante en la zona cubierta por el elemento.

\begin{ecnimportante}{Matriz de rigidez del elemento de barra de dos nodos}{cap06-elemento-2nodos-K}
Para sección transversal constante ($d\Omega = A^{(e)}\,dz$) y $E^{(e)}$ constante,
\[
\bm K^{(e)} = \int_{L^{(e)}} \bm B^{(e)\mathsf T} E^{(e)} \bm B^{(e)} \, A^{(e)}\, dz
= E^{(e)}A^{(e)} \int_0^{L^{(e)}} \begin{bmatrix} -1/L^{(e)} \\[2pt] 1/L^{(e)} \end{bmatrix} \begin{bmatrix} -1/L^{(e)} & 1/L^{(e)} \end{bmatrix} dz ,
\]
\[
\bm K^{(e)} = E^{(e)}A^{(e)} \int_0^{L^{(e)}} \begin{bmatrix} 1/(L^{(e)})^2 & -1/(L^{(e)})^2 \\ -1/(L^{(e)})^2 & 1/(L^{(e)})^2 \end{bmatrix} dz
= \frac{E^{(e)}A^{(e)}}{L^{(e)}} \begin{bmatrix} 1 & -1 \\ -1 & 1 \end{bmatrix} .
\]
\end{ecnimportante}

\begin{observacion}[Validación cruzada con el Capítulo 4]
Este resultado \emph{coincide exactamente}, término a término, con la matriz de rigidez elemental $K^{(e)}\begin{bsmallmatrix}1&-1\\-1&1\end{bsmallmatrix}$ obtenida en el Capítulo 4 mediante equilibrio directo de fuerzas internas en los extremos de la barra, con $K^{(e)}=E^{(e)}A^{(e)}/L^{(e)}$. Esta coincidencia no es casualidad: para un elemento cuya interpolación es exacta respecto del comportamiento real de la barra (deformación uniforme bajo carga axial sin distribuida), la formulación por PTV y el análisis matricial clásico deben producir --- y en efecto producen --- el mismo resultado. Esta verificación es una primera instancia, en un caso simple y comprobable a mano, del principio general de que el MEF, cuando el espacio de interpolación contiene la solución exacta del problema, la reproduce sin error (idea que se retomará formalmente en el Capítulo 7).
\end{observacion}

\begin{figure}[htbp]
\centering
\begin{tikzpicture}
\draw[->] (-0.3,0) -- (6.6,0) node[right] {$\xi = z/L^{(e)}$};
\draw[->] (0,-0.3) -- (0,1.35) node[above] {$N_i^{(e)}$};
\draw (6,0.04) -- (6,-0.04) node[below] {$1$};
\draw (0,1) -- (-0.05,1) node[left] {$1$};
\draw[ucred, thick, domain=0:6, samples=2] plot (\x,{1-\x/6});
\draw[ucblue, thick, domain=0:6, samples=2] plot (\x,{\x/6});
\filldraw (0,0) circle (1.3pt);
\filldraw (6,0) circle (1.3pt);
\node[ucred] at (0.75,1.12) {$N_1^{(e)}=1-\xi$};
\node[ucblue] at (5.4,1.12) {$N_2^{(e)}=\xi$};
\node[below] at (0,-0.35) {nodo 1};
\node[below] at (6,-0.35) {nodo 2};
\end{tikzpicture}
\caption{Funciones de forma lineales del elemento de barra de dos nodos, en la coordenada normalizada $\xi=z/L^{(e)}\in[0,1]$. Ambas rampas se cruzan en el punto medio, donde $N_1^{(e)}=N_2^{(e)}=1/2$.}
\label{fig:cap06-funciones-forma-2nodos}
\end{figure}

\section{Cargas nodales equivalentes}
\label{sec:cap06-cargas-equivalentes}

El vector de fuerzas $\bm F^{(e)}$ definido en la \cref{eq:cap06-formula-maestra} traduce cualquier carga distribuida (axial $\bar t_z$, de cuerpo $\rho b_z$) en un conjunto de fuerzas \emph{concentradas equivalentes} aplicadas en los nodos, en el sentido preciso de que producen el mismo trabajo virtual que la carga distribuida original para cualquier desplazamiento virtual cinemáticamente admisible. Esta equivalencia --- y no una intuición de "reparto de área tributaria" --- es la que gobierna el cálculo de $\bm F^{(e)}$, como se aprecia claramente en el elemento cuadrático de la sección siguiente.

Para el elemento lineal de dos nodos con $\bar t_z$, $b_z$ y $\rho$ constantes en el elemento:
\begin{align*}
\bm F^{(e)} &= \bar t_z \int_0^{L^{(e)}} \begin{bmatrix} 1-z/L^{(e)} \\ z/L^{(e)} \end{bmatrix} dz \; + \; \rho\, b_z A^{(e)} \int_0^{L^{(e)}} \begin{bmatrix} 1-z/L^{(e)} \\ z/L^{(e)} \end{bmatrix} dz \; + \; \begin{bmatrix} P_1 \\ P_2 \end{bmatrix} \\
&= \bar t_z L^{(e)} \begin{bmatrix} 1/2 \\ 1/2 \end{bmatrix} + \rho b_z A^{(e)} L^{(e)} \begin{bmatrix} 1/2 \\ 1/2 \end{bmatrix} + \begin{bmatrix} P_1 \\ P_2 \end{bmatrix} ,
\end{align*}
es decir, la resultante de la carga distribuida ($\bar t_z L^{(e)}$, o $\rho b_z A^{(e)}L^{(e)}$) se reparte por igual entre ambos nodos --- resultado intuitivo para un elemento lineal, pero que, como se verá, deja de serlo para elementos de orden superior.

Cuando la sección transversal varía linealmente entre los nodos, $A(z) = N_1^{(e)}(z)A_1 + N_2^{(e)}(z)A_2$ (es decir, se interpola con las mismas funciones de forma del desplazamiento), de modo que $d\Omega = A(z)\,dz = \left(A_1 + \dfrac{A_2-A_1}{L^{(e)}}z\right)dz$. Sustituyendo en $\bm K^{(e)} = E^{(e)}\int \bm B^{(e)\mathsf T}\bm B^{(e)}A(z)\,dz$ e integrando término a término:
\[
\bm K^{(e)} = \frac{E^{(e)}(A_1+A_2)}{2L^{(e)}} \begin{bmatrix} 1 & -1 \\ -1 & 1 \end{bmatrix},
\]
resultado formalmente idéntico al de sección constante, pero con el área reemplazada por el \emph{área promedio} $(A_1+A_2)/2$ --- una generalización natural. Para el vector de fuerzas con carga distribuida uniforme $\bar t_z$ y $\rho b_z$ constantes, la integral de $\bm N^{(e)\mathsf T}$ ponderada por $A(z)$ ya no reparte la carga de cuerpo en partes iguales, sino proporcionalmente al área asociada a cada nodo:
\[
\bm F^{(e)} = \bar t_z L^{(e)} \begin{bmatrix} 1/2 \\ 1/2 \end{bmatrix} + \rho b_z L^{(e)} \begin{bmatrix} (2A_1+A_2)/6 \\ (A_1+2A_2)/6 \end{bmatrix} + \begin{bmatrix} P_1 \\ P_2 \end{bmatrix} .
\]
Nótese que $\bar t_z$ (carga axial por unidad de longitud, ya expresada como fuerza por unidad de longitud) sigue repartiéndose por igual, mientras que $\rho b_z$ (fuerza de cuerpo por unidad de volumen) se pondera por el área local, que es mayor cerca del nodo de mayor sección.

\begin{ecnimportante}{Cargas nodales equivalentes: carga distribuida y sección variable}{cap06-cargas-equivalentes}
\textbf{Sección constante}, con $\bar t_z$, $b_z$, $\rho$ constantes en el elemento:
\[
\bm F^{(e)} = \bar t_z L^{(e)} \begin{bmatrix} 1/2 \\ 1/2 \end{bmatrix} + \rho\, b_z A^{(e)} L^{(e)} \begin{bmatrix} 1/2 \\ 1/2 \end{bmatrix} + \begin{bmatrix} P_1 \\ P_2 \end{bmatrix} .
\]
\textbf{Sección linealmente variable}, $A(z)=N_1^{(e)}A_1+N_2^{(e)}A_2$:
\begin{align*}
\bm K^{(e)} &= \frac{E^{(e)}(A_1+A_2)}{2L^{(e)}}\begin{bmatrix}1&-1\\-1&1\end{bmatrix}, \\
\bm F^{(e)} &= \bar t_z L^{(e)}\begin{bmatrix}1/2\\1/2\end{bmatrix} + \rho b_z L^{(e)} \begin{bmatrix} (2A_1+A_2)/6 \\ (A_1+2A_2)/6 \end{bmatrix} + \begin{bmatrix}P_1\\P_2\end{bmatrix} .
\end{align*}
En ambos casos, la carga nodal equivalente se obtiene por integración exacta de $\bm N^{(e)\mathsf T}$ contra la carga distribuida y \emph{no} por un reparto geométrico intuitivo de área tributaria (aunque, para estos casos particulares, ambos criterios coinciden).
\end{ecnimportante}

\begin{figure}[htbp]
\centering
\begin{tikzpicture}
\draw[thick, fill=ucblue!8, draw=ucblue] (0,-0.9) -- (6,-0.4) -- (6,0.4) -- (0,0.9) -- cycle;
\filldraw (0,0) circle (2pt);
\filldraw (6,0) circle (2pt);
\node[below=16pt] at (0,0) {nodo 1, $A_1$};
\node[below=16pt] at (6,0) {nodo 2, $A_2$};
\foreach \x in {0.5,1.5,2.5,3.5,4.5,5.5} {
  \pgfmathsetmacro{\yy}{0.4+0.5*(\x/6)}
  \draw[->, ucred, thick] (\x,{\yy+0.55}) -- (\x,{\yy+0.05});
}
\node[above, ucred] at (3,1.65) {$\bar t_z(z)$};
\draw[->] (6.4,0) -- (7.1,0) node[right] {$z$};
\end{tikzpicture}
\caption{Barra de sección transversal linealmente variable entre $A_1$ y $A_2$, sometida a una carga axial distribuida $\bar t_z(z)$.}
\label{fig:cap06-barra-seccion-variable}
\end{figure}

\section{Elemento cuadrático de tres nodos}
\label{sec:cap06-elemento-3nodos}

Cuando se dispone de un tercer nodo (nodo 3, ubicado --- por simplicidad --- en el punto medio del elemento, a una distancia $L^{(e)}/2$ del nodo 1), la interpolación puede enriquecerse a un polinomio \emph{cuadrático}:
\[
w = a_0 + a_1 z + a_2 z^2 .
\]
Imponiendo los tres valores nodales,
\begin{align*}
w\big|_{z=0} &=W_1^{(e)}=a_0, \\
w\big|_{z=L^{(e)}} &=W_2^{(e)}=a_0+a_1L^{(e)}+a_2(L^{(e)})^2, \\
w\big|_{z=L^{(e)}/2} &=W_3^{(e)}=a_0+a_1\frac{L^{(e)}}{2}+a_2\frac{(L^{(e)})^2}{4},
\end{align*}
se obtiene un sistema lineal de tres ecuaciones para $a_0,a_1,a_2$ cuya solución es
\[
a_0=W_1^{(e)}, \qquad a_1=\frac{-3W_1^{(e)}-W_2^{(e)}+4W_3^{(e)}}{L^{(e)}}, \qquad a_2=\frac{2W_1^{(e)}+2W_2^{(e)}-4W_3^{(e)}}{(L^{(e)})^2} .
\]
Sustituyendo de vuelta y reagrupando en términos de $z/L^{(e)}$, se identifican las tres funciones de forma cuadráticas:
\[
N_1^{(e)} = 1 - 3\frac{z}{L^{(e)}} + 2\left(\frac{z}{L^{(e)}}\right)^2, \qquad
N_2^{(e)} = -\frac{z}{L^{(e)}} + 2\left(\frac{z}{L^{(e)}}\right)^2, \qquad
N_3^{(e)} = 4\frac{z}{L^{(e)}} - 4\left(\frac{z}{L^{(e)}}\right)^2 .
\]
Puede verificarse directamente que $N_i^{(e)}(z_j)=\delta_{ij}$ para los tres nodos, y que $N_1^{(e)}+N_2^{(e)}+N_3^{(e)}=1$ idénticamente. Nótese, sin embargo, que $N_1^{(e)}$ y $N_2^{(e)}$ no son monótonas ni positivas en todo el elemento: cada una alcanza un valor levemente negativo (mínimo $-1/8$) en la mitad del elemento más alejada de su propio nodo, mientras que $N_3^{(e)}$, asociada al nodo central, describe una parábola simple con máximo unitario en $z=L^{(e)}/2$ (Figura~\ref{fig:cap06-funciones-forma-3nodos}).

La matriz de deformación se obtiene derivando cada función de forma:
\[
\bm B^{(e)} = \begin{bmatrix} -\dfrac{3}{L^{(e)}}+\dfrac{4z}{(L^{(e)})^2} & -\dfrac{1}{L^{(e)}}+\dfrac{4z}{(L^{(e)})^2} & \dfrac{4}{L^{(e)}}-\dfrac{8z}{(L^{(e)})^2} \end{bmatrix} ,
\]
que ahora es \emph{lineal} en $z$ (a diferencia del caso de dos nodos, donde era constante): el elemento cuadrático puede representar un estado de deformación linealmente variable dentro de su dominio, una mejora significativa de la capacidad de aproximación.

\begin{ecnimportante}{Matriz de rigidez del elemento de barra de tres nodos}{cap06-elemento-3nodos-K}
Con $E^{(e)}$, $A^{(e)}$ constantes, $\bm K^{(e)}=\displaystyle\int_0^{L^{(e)}} \bm B^{(e)\mathsf T}E^{(e)}\bm B^{(e)}A^{(e)}\,dz$ se calcula integrando término a término el producto de polinomios de $\bm B^{(e)\mathsf T}\bm B^{(e)}$ (todos de grado 2 en $z$), obteniéndose el resultado cerrado
\[
\bm K^{(e)} = \frac{E^{(e)}A^{(e)}}{L^{(e)}} \begin{bmatrix} 7/3 & 1/3 & -8/3 \\ 1/3 & 7/3 & -8/3 \\ -8/3 & -8/3 & 16/3 \end{bmatrix} .
\]
Como verificación de consistencia, obsérvese que cada fila (y, por simetría, cada columna) suma cero: $7/3+1/3-8/3=0$ y $-8/3-8/3+16/3=0$. Esto refleja que, ante un desplazamiento de sólido rígido aplicado a los tres nodos ($W_1^{(e)}=W_2^{(e)}=W_3^{(e)}$), las fuerzas internas resultantes $\bm K^{(e)}\bm U^{(e)}$ deben autoequilibrarse (suma nula), tal como se comprobó para el elemento de dos nodos en el Capítulo 4.
\end{ecnimportante}

Para las cargas nodales equivalentes, con $\bar t_z$, $b_z$ y $\rho$ constantes en el elemento, la integración de $\bm N^{(e)\mathsf T}$ (que ahora incluye las tres funciones de forma cuadráticas) entrega:

\begin{ecnimportante}{Lumping de carga uniforme para el elemento cuadrático}{cap06-lumping-cuadratico}
\[
\bm F^{(e)} = \bar t_z L^{(e)} \begin{bmatrix} 1/6 \\ 1/6 \\ 4/6 \end{bmatrix} + \rho\, b_z A^{(e)} L^{(e)} \begin{bmatrix} 1/6 \\ 1/6 \\ 4/6 \end{bmatrix} + \begin{bmatrix} P_1 \\ P_2 \\ P_3 \end{bmatrix} .
\]
Nótese el resultado, en apariencia contraintuitivo: aunque la carga distribuida es \emph{uniforme} a lo largo del elemento, su reparto nodal equivalente \emph{no} es uniforme ($1/6, 1/6, 4/6$ en vez de $1/3,1/3,1/3$). El nodo central recibe dos tercios de la resultante total, y los nodos extremos solo un sexto cada uno. Este resultado es consecuencia directa de la forma de las funciones de forma cuadráticas (en particular, de que $N_1^{(e)}$ y $N_2^{(e)}$ toman valores negativos en parte del dominio) y no admite una interpretación simple en términos de "longitud tributaria"; es, precisamente, la equivalencia en el sentido del PTV la que gobierna el resultado.
\end{ecnimportante}

\begin{figure}[htbp]
\centering
\begin{tikzpicture}
\draw[->] (-0.3,0) -- (6.6,0) node[right] {$\xi=z/L^{(e)}$};
\draw[->] (0,-0.35) -- (0,1.3) node[above] {$N_i^{(e)}$};
\draw (6,0.03) -- (6,-0.03) node[below] {$1$};
\draw (3,0.03) -- (3,-0.03) node[below] {$1/2$};
\draw (0,1) -- (-0.05,1) node[left] {$1$};
\draw[dashed, gray!60] (-0.3,0) -- (6.6,0);
\draw[ucred, thick, domain=0:6, samples=80] plot (\x,{1-3*(\x/6)+2*(\x/6)*(\x/6)});
\draw[ucblue, thick, domain=0:6, samples=80] plot (\x,{-(\x/6)+2*(\x/6)*(\x/6)});
\draw[ucteal, thick, domain=0:6, samples=80] plot (\x,{4*(\x/6)-4*(\x/6)*(\x/6)});
\filldraw (0,0) circle (1.3pt);
\filldraw (3,0) circle (1.3pt);
\filldraw (6,0) circle (1.3pt);
\node[ucred] at (0.9,1.12) {$N_1^{(e)}$};
\node[ucblue] at (5.6,1.12) {$N_2^{(e)}$};
\node[ucteal] at (3,1.12) {$N_3^{(e)}$};
\end{tikzpicture}
\caption{Funciones de forma cuadráticas del elemento de barra de tres nodos. $N_1^{(e)}$ y $N_2^{(e)}$ alcanzan un mínimo levemente negativo ($-1/8$) en la mitad del elemento opuesta a su nodo; $N_3^{(e)}$ (nodo central) es una parábola positiva de máximo unitario.}
\label{fig:cap06-funciones-forma-3nodos}
\end{figure}

\begin{figure}[htbp]
\centering
\begin{subfigure}[b]{0.46\textwidth}
\centering
\begin{tikzpicture}
\draw[thick] (0,0) -- (4,0);
\filldraw (0,0) circle (2pt) node[below=4pt] {$1$};
\filldraw (4,0) circle (2pt) node[below=4pt] {$2$};
\draw[->, ucred, thick] (0,0) -- (0,0.8) node[above] {$W_1^{(e)}$};
\draw[->, ucred, thick] (4,0) -- (4,0.8) node[above] {$W_2^{(e)}$};
\draw[->] (4.3,0) -- (5,0) node[right] {$z$};
\end{tikzpicture}
\caption{Elemento de dos nodos}
\label{fig:cap06-gdl-2nodos}
\end{subfigure}
\hfill
\begin{subfigure}[b]{0.46\textwidth}
\centering
\begin{tikzpicture}
\draw[thick] (0,0) -- (4,0);
\filldraw (0,0) circle (2pt) node[below=4pt] {$1$};
\filldraw (2,0) circle (2pt) node[below=4pt] {$3$};
\filldraw (4,0) circle (2pt) node[below=4pt] {$2$};
\draw[->, ucred, thick] (0,0) -- (0,0.8) node[above] {$W_1^{(e)}$};
\draw[->, ucred, thick] (2,0) -- (2,0.8) node[above] {$W_3^{(e)}$};
\draw[->, ucred, thick] (4,0) -- (4,0.8) node[above] {$W_2^{(e)}$};
\draw[->] (4.3,0) -- (5,0) node[right] {$z$};
\end{tikzpicture}
\caption{Elemento de tres nodos (nodo 3 en el punto medio)}
\label{fig:cap06-gdl-3nodos}
\end{subfigure}
\caption{Grados de libertad nodales (desplazamientos axiales) de los elementos de barra de dos y de tres nodos.}
\label{fig:cap06-gdl-elementos}
\end{figure}

\section{Elementos de barra Lagrangianos generales}
\label{sec:cap06-lagrangianos}

Los casos de dos y tres nodos desarrollados anteriormente son instancias particulares de una construcción general: el uso de \emph{polinomios de interpolación de Lagrange} para definir las funciones de forma de un elemento de barra con un número arbitrario $n$ de nodos, todos ellos de clase $C^0$. Un polinomio de Lagrange asociado al nodo $i$ se construye exigiendo que valga $1$ en $z_i$ y $0$ en los restantes $n-1$ nodos $z_j$ ($j\neq i$):

\begin{ecnimportante}{Funciones de forma Lagrangianas generales}{cap06-lagrangianos-generales}
Para un elemento de barra de $n$ nodos ubicados en las posiciones $z_i = \alpha_i L^{(e)}$ ($i=1,\dots,n$, con $\alpha_i\in[0,1]$ los parámetros de posición relativa),
\[
N_i^{(e)}(z) = \prod_{\substack{j=1 \\ j\neq i}}^{n} \frac{z-z_j}{z_i-z_j} = \prod_{\substack{j=1\\j\neq i}}^{n} \frac{z/L^{(e)}-\alpha_j}{\alpha_i-\alpha_j}, \qquad i=1,\dots,n .
\]
Por construcción, $N_i^{(e)}(z_j)=\delta_{ij}$, y el grado del polinomio resultante es $n-1$ (un polinomio de grado $n-1$ queda determinado de forma única por sus $n$ valores nodales).
\end{ecnimportante}

Como verificación, para $n=2$ con $\alpha_1=0,\ \alpha_2=1$:
\[
N_1^{(e)} = \frac{z/L^{(e)}-1}{0-1} = 1-\frac{z}{L^{(e)}}, \qquad N_2^{(e)}=\frac{z/L^{(e)}-0}{1-0}=\frac{z}{L^{(e)}} ,
\]
que reproduce exactamente el resultado de la \cref{sec:cap06-elemento-2nodos}. Para $n=3$ con $\alpha_1=0,\ \alpha_2=1,\ \alpha_3=1/2$:
\[
N_1^{(e)} = \frac{\left(z/L^{(e)}-1\right)\left(z/L^{(e)}-1/2\right)}{(0-1)(0-1/2)} = 1-3\frac{z}{L^{(e)}}+2\left(\frac{z}{L^{(e)}}\right)^2 ,
\]
e, idénticamente, $N_2^{(e)}$ y $N_3^{(e)}$ reproducen los resultados de la \cref{sec:cap06-elemento-3nodos}. Esta construcción general es la que permite --- como se explotará en el Capítulo 8 --- definir elementos de orden arbitrariamente alto sin necesidad de resolver un sistema de ecuaciones distinto para cada valor de $n$, y es también la base de la formulación isoparamétrica, en la que las mismas funciones de Lagrange se emplean para interpolar tanto el campo de desplazamientos como la geometría del elemento.

\section{Síntesis}

En este capítulo se completó el programa anunciado al cierre del Capítulo 5: la construcción sistemática de las matrices elementales $\bm K^{(e)}$ y $\bm F^{(e)}$ a partir de tres ingredientes --- interpolación $\bm u=\bm N\bm U$, deformación $\bm\varepsilon=\bm B\bm U$ y constitutiva $\bm\sigma=\bm C\bm\varepsilon$ --- sustituidos en el PTV. Se verificó, para el elemento lineal de dos nodos, la coincidencia exacta con el resultado del análisis matricial clásico (Capítulo 4), y se extendió la formulación al elemento cuadrático de tres nodos y, finalmente, a una familia Lagrangiana de orden arbitrario. Quedan pendientes dos preguntas centrales que se abordan en los capítulos siguientes: ¿bajo qué condiciones garantiza esta construcción que la solución numérica converja a la solución exacta al refinar la malla? (Capítulo 7); y ¿cómo se pueden evaluar de manera sistemática y eficiente las integrales que definen $\bm K^{(e)}$ y $\bm F^{(e)}$, especialmente cuando la geometría del elemento no es tan simple como la de una barra recta? (Capítulo 8).
